\chapter{Numerical issues}
\label{chap:num}

\section{Choice probability computation}

We start from the logit expression
\begin{equation}
P_{ij} = \frac{e^{U_{\ell}}}{\sum_{k \in \mathcal{A}_i} e^{U_k}}.
\label{eq:probab}
\end{equation}
If some of the utilities are too large, the computation of
\eqref{eq:probab} can face some overflow.
A simple way to avoid this problem is to substract some quantity
$\alpha$ since
\[
\frac{e^{U_{\ell}-\alpha}}{\sum_{k \in \mathcal{A}_i} e^{U_k-\alpha}}
= \frac{e^{U_{\ell}}}{\sum_{k \in \mathcal{A}_i} e^{U_k}}.
\] 
Moreover, is all utilities are too small, they could be evaluated to
0, leading to the ratio $0/0$ in \eqref{eq:probab}.
A simple way to circumvent this issue is to normalize with respect to
the highest utility, i.e. take
\[
\alpha = \max \lbrace U_k,\ k \in \mathcal{A}_i \rbrace - \beta.
\]
We want to keep each term $e^{U_k}$, $k \in \mathcal{A}_i$, in the
same range, so $\beta$ has to be kept low.
Since $d e^x/dx = e^x$, the precise value of $\beta$ does not
significantly affect small utilities, so a simple, while practical,
choice is to set $\beta$ to 0.
As a consequence, $e^{U_k - \alpha} \in (0,1]$ for all $k \in \mathcal{A}_i$.
